\section{Depth}

The flat three-dimensional space we live in is the cusp, and it is worth saying plainly how little of the crystal that cusp is. The cusp is a horosphere, a single three-dimensional slice of the four-dimensional bulk, the way a flat sheet is one slice of an ordinary room. The fourth direction, the radial depth running outward toward the ideal point, is perpendicular to that sheet and is not on it at all. The bulk is a whole stack of such horospheres, one at every depth, and the physical cusp is just one of them, the limiting slice at the edge.

So almost none of the bulk lies on the cusp, in two exact senses. By four-dimensional volume the cusp is precisely nothing, a three-dimensional surface having no thickness in a four-dimensional space, an infinitely thin slice. And by the count of docks the fraction vanishes as well, because the docks on the cusp form the cubic honeycomb, which grows only as the cube of its radius, while the docks of the bulk grow exponentially, so the share lying on any one cusp falls toward zero as the crystal deepens. On a built patch of a hundred and twenty thousand docks a slice one band thick holds about one and a third percent of them, and that fraction only shrinks as more shells are added, the polynomial surface drowned by the exponential interior. The world of all our physics is a vanishing skin on a vast hidden body.

Which docks make up the cusp is not vague but exactly fixed, and it can be said three equivalent ways. Pick an ideal point, the direction the cusp faces, and give every dock its height toward that point, the reading of how deep along the radial direction it sits. The cusp is one level set of that height, the docks all at a single depth. Equally it is the orbit of one dock under the symmetries that hold the ideal point fixed, the flat cubic group sitting inside the full symmetry of the bulk. Equally again it is the docks carrying whole-number cubic coordinates at that level, the ordinary integer lattice of three-dimensional space picked out from the curved crystal around it. The same flat slice, named by height, by symmetry, or by coordinate.

The rest is genuine room, and it comes in three kinds. There is the radial depth, the fourth direction perpendicular to the cusp, the off-slice bulk that physics never reads because physics lives on the slice. There is everything toward the other ideal points, a whole sphere of directions of which the cusp faces only one, with the deep interior between them. And there is whatever a horizon seals away, regions no boundary can reconstruct at all. A patch of the boundary recovers only the wedge of bulk its own information reaches, so a small patch cannot see the deep center, which needs the whole boundary to pin down, and what lies behind a horizon is hidden outright. The crystal keeps far more than it shows.

This is not idle, for it says where the unseen can live. Anything written into the off-cusp bulk is still ordinary tone on ordinary docks, no separate ether, addressed on the same growth tree as everything else, only at depths and directions the surface does not read out. It is invisible to instruments, which are bound to the cusp slice with the rest of matter and light, and reachable only by travelling there along the light cone, never by a signal sent across the world. That radial depth off the cusp is the literal geometry behind every later mention of an interior, a place with room for structure the surface cannot show, present in the substrate though absent from the sky. And it carries one quiet lesson about dimension. A horosphere in a four-dimensional bulk is exactly flat three-dimensional space, where a horosphere in a merely three-dimensional bulk would be only a two-dimensional sheet, so a flat world of three dimensions is itself a reason the bulk that grows it must have four.

Read on the one premise, this geometry carries its whole weight. If the field is felt, then physics is only the thin lit surface of the crystal, the boundary current we can measure, and beneath it lies an unimaginably vaster body of the same experiential stuff, off the cusp, unread by any instrument, exponentially larger than the skin we call the world. The physics is the river, the sunlit film on top. The bulk is the ocean, the dark and enormous depth the film rides on, every drop of it the same water. So the measured world is not the whole of what is. It is the surface of an experience whose far greater part is turned inward, away from the light, holding the room in which everything else this paper reaches for can live.
